\documentclass[letterpaper,10pt,conference]{ieeeconf}
\IEEEoverridecommandlockouts
\overrideIEEEmargins

\usepackage{amsmath,amssymb,bm}
\usepackage{booktabs}
\usepackage{graphicx}
\usepackage{multirow}
\usepackage{url}
\usepackage{cite}
\usepackage{xcolor}
\usepackage{tikz}
\usetikzlibrary{arrows.meta,positioning,fit}

% Citation/link blue sampled from the CoSAM manuscript (#367DBD).
\definecolor{cosamLinkBlue}{HTML}{367DBD}

% ieeeconf defines an obsolete Natbib compatibility hook.  Removing it before
% loading hyperref restores clickable numerical citations.
\makeatletter
\let\NAT@parse\undefined
\makeatother
\usepackage[
  colorlinks=true,
  linkcolor=cosamLinkBlue,
  citecolor=cosamLinkBlue,
  urlcolor=black
]{hyperref}

\definecolor{todoRed}{RGB}{176,36,36}
\definecolor{inputBlue}{RGB}{225,238,250}
\definecolor{geometryGreen}{RGB}{224,243,231}
\definecolor{routeOrange}{RGB}{253,237,214}
\definecolor{outputPurple}{RGB}{238,229,247}
\definecolor{transportTeal}{HTML}{12857D}
\definecolor{unsupportedPurple}{HTML}{7455A4}
\newcommand{\todo}[1]{\textcolor{todoRed}{\textbf{[TODO: #1]}}}
\newcommand{\tbd}{\textcolor{todoRed}{--}}
\newcommand{\captioncolorchip}[2]{%
  \mbox{(\raisebox{0.05ex}{\textcolor{#1}{\rule{0.9ex}{0.9ex}}}%
  \hspace{0.35em}\textcolor{#1}{#2})}}
\newcommand{\R}{\mathbb{R}}
\newcommand{\I}{\mathbf{I}}
\newcommand{\D}{\mathbf{D}}
\newcommand{\K}{\mathbf{K}}
\newcommand{\T}{\mathbf{T}}
\newcommand{\M}{\mathbf{M}}
\newcommand{\F}{\mathbf{F}}
\newcommand{\U}{\mathbf{U}}
\newcommand{\n}{\mathbf{n}}
\newcommand{\od}{\mathbin{\odot}}
\newsavebox{\roboticplatformbox}

\title{\LARGE \bf
DistriNVS: Distributional Geometry Transport and Visibility-Constrained Completion for Cross-Endoscope Novel-View Synthesis}

\author{Anonymous Authors}

\begin{document}
\maketitle
\thispagestyle{empty}
\pagestyle{empty}

\begin{abstract}

Cross-endoscope novel-view synthesis must transport anatomy observed by the
source camera while synthesizing tissue revealed only from the target
viewpoint. These regimes become difficult to separate when sparse or erroneous
depth is projected as deterministic points, because sampling gaps, depth
failures, surface collisions, and true disocclusions can produce
indistinguishable holes. We present DistriNVS, a reliability-aware framework
that represents source depth as a distribution and propagates it through
calibrated geometry. Distributional surface splatting estimates target support,
depth variance, and collision ambiguity, which guide a soft mixture between a
bounded transport expert and an uncertainty-conditioned Fourier synthesis
expert.
DistriNVS is trained on 13 iMED scenes, evaluated on 1,392 iMED test frames,
and evaluated zero-shot on 275 EndoVis test pairs. On the iMED test set,
DistriNVS achieves a PSNR of 22.95~dB, an SSIM of 0.59, and an H-PSNR of
27.95~dB. In zero-shot evaluation on EndoVis, DistriNVS achieves the best
PSNR of 23.47~dB, SSIM of 0.66, and LPIPS of 0.23 among the compared
feed-forward methods, together with an H-PSNR of 21.11~dB and an inference
time of 0.08~s per frame. These results demonstrate the effectiveness of
reliability-guided transport--synthesis fusion for cross-endoscope novel-view
synthesis and its generalization to an unseen dataset without target-domain
adaptation.
\textit{\textcolor{cosamLinkBlue}{The code will be made publicly available upon publication.}}

\end{abstract}

\section{Introduction}
\label{sec:introduction}

Robot-assisted minimally invasive surgery depends on endoscopic vision for
scene perception and instrument guidance
~\cite{wei2026roboticperception,sunmola2026surgicalgaussiansurfels}.
Cross-endoscope novel-view synthesis can broaden this visual context by
predicting a target-endoscope view from synchronized appearance and geometry
acquired by a spatially displaced source endoscope~\cite{wang2026endoe2egs}.
Unlike left-to-right synthesis within one stereo endoscope, the target may
observe anatomy absent from both source cameras. The task therefore combines
two distinct regimes: co-visible anatomy has measured appearance
that should be transported through calibrated geometry, whereas target-only
anatomy lacks source evidence and must be synthesized. A reliable system should
retain source evidence, identify inferred content, and thereby support
wider-field robotic perception, retrospective visualization, and surgical
simulation~\cite{hsieh2026extrags}.
Such uncertainty-aware cross-view completion may also support visual-state
construction for surgical robot world models~\cite{he2025surgworld}.

Achieving this separation is difficult because endoscopic depth and calibration
are unreliable around specularities, weak texture, depth discontinuities,
instruments, and deforming tissue~\cite{allan2021scared}. Conventional forward
warping represents each valid source pixel by one 3-D point; discretized
projection leaves sampling gaps, while multiple points may compete for the same
target pixel. A binary validity mask then conflates missing depth, rasterization
holes, depth-ordering conflicts, and genuine cross-endoscope disocclusions.
Treating these cases uniformly either degrades observed structure or permits
hallucination where geometric evidence remains trustworthy.

\begin{figure}[t]
\centering
\vspace*{2mm}
\sbox{\roboticplatformbox}{%
  \includegraphics[width=\columnwidth]{figures/distrinvs_teaser/distrinvs_teaser_editable.pdf}%
}
\begin{tikzpicture}
  \clip[rounded corners=0mm]
    (0,0) rectangle
    (\wd\roboticplatformbox,\ht\roboticplatformbox);
  \node[anchor=south west,inner sep=0pt] at (0,0)
    {\usebox{\roboticplatformbox}};
\end{tikzpicture}
\caption{Cross-endoscope novel-view synthesis under partial field-of-view
overlap. Given the source-to-target pose, co-visible anatomy is
source-supported and can be transported geometrically
\captioncolorchip{transportTeal}{teal}, whereas target-only anatomy lacks source
evidence and needs to be synthesized \captioncolorchip{unsupportedPurple}{purple}.
}
\label{fig:conditional_refinement}
\vspace{-0.4cm}
\end{figure}

Existing methods address complementary parts of this problem but do not expose
the required evidence boundary. Differentiable point rendering and softmax
splatting improve feature transport and collision handling
~\cite{wiles2020synsin,niklaus2020softmax}, while reference-guided inpainting
combines depth-based alignment with learned completion~\cite{zhao2023geofill}.
Large-mask inpainting networks provide strong global image priors
~\cite{suvorov2022lama,chu2023uffc}, and endoscopic video completion adds
depth, temporal, or stereo cues
~\cite{daher2023temporal,zhang2024daevi,zou2025ssifnet}. Endoscopic neural
rendering has likewise advanced sparse-view reconstruction and uncertainty
modeling~\cite{li2024endosparse,guo2025ucnerf}. However, deterministic
projection assigns equal reliability to all valid depths, while image-space
uncertainty does not describe how source geometry
supports a spatially offset target view. Neither approach propagates per-pixel
depth uncertainty through the cross-endoscope transform and uses the resulting
visibility statistics to constrain where appearance may be synthesized.

Our central premise is to \emph{transport reliable geometric evidence and
synthesize only unsupported or geometrically unreliable regions}. We introduce
DistriNVS, a reliability-aware framework
for direct cross-endoscope novel-view synthesis. Its distributional surface
splatting (DSS) represents each source depth as a distribution whose bounded
mean and scale, surface normal, and confidence are predicted by a depth
reliability encoder. DSS projects deterministic sigma samples through the
calibrated cross-endoscope transform and aggregates source features using
Jacobian- and normal-aware footprints with soft nearest-surface visibility,
yielding target support, depth variance, and collision entropy. These statistics
constrain a soft mixture between bounded residual transport for co-visible
regions and uncertainty-conditioned Fourier completion for unsupported or
ambiguous regions. Target-depth and risk prediction plus backward RGB-D
reprojection constrain this mixture during training.

The paper makes three contributions:
\begin{itemize}
    \item We propose a reliability-aware framework for cross-endoscope novel-view 
    synthesis that separates the transport of source-observed anatomy from the 
    synthesis of source-unobserved content.
    \item We introduce DSS, a differentiable renderer that propagates depth
    distributions through calibrated cross-endoscope geometry and exposes
    target support, depth variance, and collision entropy beyond binary validity.
    \item We design a two-expert decoder that forms a soft mixture of bounded
    residual transport and uncertainty-conditioned Fourier completion under a
    renderer-derived reliability constraint.
\end{itemize}

\section{Related Work}
\label{sec:related}

We review image completion, geometry-guided view transport, and
endoscopic neural rendering, with particular attention to how existing
approaches represent missing content, exploit calibrated geometry, and quantify
uncertainty.

\begin{figure*}[t]
\centering
\includegraphics[width=\textwidth]{figures/distrinvs_framework/distrinvs_framework.pdf}
\caption{DistriNVS framework. A reliability encoder predicts source geometry
and features. DSS transports deterministic depth hypotheses through the
calibrated cross-endoscope transform and jointly produces warped evidence and
target-plane visibility statistics. These statistics condition a
reliability-constrained soft mixture of a residual transport expert and a
Fourier synthesis expert. Solid arrows denote inference; the dashed lower lane is used only
during training.}
\label{fig:overview}
\vspace{-0.4cm}
\end{figure*}

\subsection{Image and Endoscopic Video Inpainting}

Generic image inpainting combines mask-aware feature extraction with
broad appearance priors. Partial and gated convolutions prevent
invalid activations from contaminating local features
~\cite{liu2018partial,yu2019gated}, whereas LaMa and UFFC use Fourier operations
for global context and MAT restricts attention to valid tokens
~\cite{suvorov2022lama,chu2023uffc,li2022mat}. Recent diffusion models such as
BrushNet and TurboFill incorporate masked-image features into pretrained
generators or adapt few-step sampling for efficient completion
~\cite{ju2024brushnet,xie2025turbofill}. Despite strong appearance priors, these
methods define holes in a fixed image coordinate system and do not describe
geometric provenance. Endoscopic completion adds domain cues: temporal learning
removes specularities, DAEVI combines spatio-temporal RGB and depth, and SSIFNet
fuses flow and stereo-temporal information for instrument occlusions
~\cite{daher2023temporal,zhang2024daevi,zou2025ssifnet}. Such methods exploit
neighboring frames or views but still restore defects defined in the original
camera coordinates. Cross-endoscope synthesis instead creates holes through
viewpoint transfer: target-only anatomy is unobserved, while rasterization gaps
or depth conflicts may retain source support. DistriNVS therefore uses
calibrated projection statistics to distinguish true disocclusions from
rasterization gaps and depth conflicts.

\subsection{Geometry-Guided Transport and Completion}

Geometry-guided view synthesis preserves correspondence by transporting source
evidence before generating missing content. SynSin predicts and projects a
feature point cloud, softmax splatting resolves many-to-one mappings with
learned importance, and GeoFill refines depth and pose before reprojection and
inpainting~\cite{wiles2020synsin,niklaus2020softmax,zhao2023geofill}. Recent
feed-forward Gaussian renderers use stronger explicit representations: MVSplat
localizes Gaussians with plane-sweep cost volumes, whereas DepthSplat combines
multi-view correspondence with monocular depth priors
~\cite{chen2024mvsplat,xu2025depthsplat}. They reduce per-scene optimization but
infer a radiance representation from multiple calibrated observations, without
an explicit per-pixel boundary between measured and synthesized appearance.
Generative methods target wider gaps. LeftRefill frames reference-conditioned
synthesis as paired-canvas completion, MultiDiff conditions video diffusion on
depth-warped reference views, and point-guided latent diffusion uses
stereo-derived clouds for laparoscopic view synthesis
~\cite{cao2024leftrefill,muller2024multidiff,tang2025pointguided}. These models can
generate plausible structures in large disocclusions, yet plausibility alone
does not reveal which pixels remain geometrically supported. DistriNVS instead
transports a calibrated source observation with distributional splatting,
derives support, variance, and collision statistics during rasterization, and
uses them to constrain soft fusion between transported and synthesized content.

\subsection{Endoscopic Novel-View Synthesis and Uncertainty}

Endoscopic neural rendering must handle non-rigid tissue, view-dependent
illumination, tool occlusions, and sparse camera coverage. EndoNeRF represents
deformation with a neural radiance field, while EndoSurf combines dynamic
modeling with a neural surface~\cite{wang2022endonerf,zha2023endosurf}.
EndoSparse introduces geometric and appearance priors into Gaussian splatting,
and Endo-GSMT uses depth, dense-displacement, and motion priors for dynamic
reconstruction~\cite{li2024endosparse,gou2025endogsmt}. GastroNVS further shows
that large viewpoint changes in real gastroscopy amplify geometric artifacts
and inter-view illumination mismatch~\cite{minai2026gastronvs}. These systems
build persistent scene representations from multiple frames and address
reconstruction fidelity within a captured scene. Uncertainty-aware
approaches address sparse evidence from another direction: UC-NeRF conditions a
radiance field on multi-view uncertainty, whereas ExtraGS uses uncertainty to
sample virtual cameras and weight diffusion-refined pseudo-observations for
extrapolation~\cite{guo2025ucnerf,hsieh2026extrags}. Thus, neither uncertainty nor
generative priors are new to endoscopic rendering. The distinction is
operational: DSS propagates a distribution for each input depth through the
calibrated cross-endoscope transform and derives target support, depth variance,
and collision entropy from the same rasterization that transports appearance.
DistriNVS uses these quantities to route completion during unseen-sequence
inference without target-scene optimization.

\section{Method}
\label{sec:method}

% \subsection{Problem Formulation}

% After resizing, both endoscope views have spatial size $H\times W$; we omit the
% batch dimension throughout. 
DistriNVS receives source RGB
$\I_s\in[0,1]^{3\times H\times W}$, metric depth
$\D_s\in\R_+^{1\times H\times W}$, and validity mask
$\M_s\in\{0,1\}^{1\times H\times W}$. Calibration comprises intrinsics
$\K_s,\K_t\in\R^{3\times3}$ and a rigid transform
$\T_{s\rightarrow t}\in\mathrm{SE}(3)\subset\R^{4\times4}$. The model predicts
target RGB $\widehat\I_t\in[0,1]^{3\times H\times W}$, depth
$\widehat\D_t\in\R_+^{1\times H\times W}$, and risk
$\U_t\in[0,1]^{1\times H\times W}$. 
% Target RGB-D and the target instrument mask
% are used only for training or evaluation.

\subsection{Reliability Encoder}

Fig.~\ref{fig:overview} summarizes the inference pipeline. A
reliability encoder estimates source geometry and uncertainty. DSS then transports source evidence and measures target-view
support and ambiguity. These measurements route complementary transport and
synthesis experts, whose RGB-D predictions are combined by a
reliability-constrained soft gate. The backward photometric cycle and risk objective
are used only during training.
We first construct the source-side geometric representation on which DSS
operates. Rather than treating measured and inferred depths as equally
reliable, the reliability encoder processes source RGB, normalized sparse
depth, and its validity mask. A full-resolution branch preserves local
boundaries through residual convolutional blocks, while a quarter-resolution
branch captures broader context using strided convolution. The context features
are bilinearly upsampled, concatenated with the full-resolution features, and
residually fused. A shared prediction head produces a feature map
$\F_s\in\R^{C_f\times H\times W}$ and raw depth, scale, confidence, and
orientation parameters. These parameters yield depth fields
$\mu,\sigma\in\R_+^{1\times H\times W}$, confidence
$c\in[0,1]^{1\times H\times W}$, and a unit-length orientation proxy
$\n\in\R^{3\times H\times W}$.

Measured depths receive bounded residual corrections, whereas missing depths
are inferred through a \emph{depth-completion (DC)} parameterization of the
shared prediction head. When DC is enabled, both measured and completed depths
participate in DSS, with reduced confidence assigned to completed estimates.
The orientation proxy is learned through downstream objectives and is used
to adapt the splat footprint. During training, random depth dropout exposes the
encoder to both measured and missing geometry while the original valid samples
retain depth supervision.

\subsection{Distributional Surface Splatting}

Given these source-side estimates, DSS transports the source attributes
$\mathbf A_s=[\I_s,\F_s]\in\R^{(3+C_f)\times H\times W}$ to
$\mathbf A_t^w\in\R^{(3+C_f)\times H\times W}$ and estimates expected target
depth $\bar D_t\in\R_+^{1\times H\times W}$. The same rasterization produces
support and depth variance $S,V\in\R_+^{1\times H\times W}$; coverage,
collision entropy, and confidence
$C,\mathcal H,\kappa\in[0,1]^{1\times H\times W}$; and render validity
$M^{\rm dss}\in\{0,1\}^{1\times H\times W}$.
These maps indicate which target pixels are 
supported by source geometry and how reliable that support is. We define
$\Omega_s=\Omega_t=\{1,\ldots,W\}\times\{1,\ldots,H\}
\subset\mathbb Z^2$ as the source and target pixel grids. 

\textbf{Distributional depth (DD).}
A source pixel
$q\in\Omega_s$ and index $k\in\mathcal K=\{1,2,3\}$ identify a depth
hypothesis. We place these hypotheses
at the predicted mean and one standard deviation on either side, with
normalized Gaussian weights $a_k\in[0,1]$ satisfying
$\sum_{k\in\mathcal K}a_k=1$. The calibrated pinhole operator projects each
hypothesis $z_{qk}\in\R_+$ to a target coordinate
$\mathbf x_{qk}\in\R^2$ and target depth $z^t_{qk}\in\R_+$. It comprises source
back-projection, rigid transformation, and target projection
~\cite{wiles2020synsin,niklaus2020softmax}.

\textbf{Orientation-aware footprints (OF).}
Depth uncertainty is converted into a hypothesis-dependent, axis-aligned
target footprint. The projection derivative $J^\alpha_{qk}$ measures the
sensitivity of image coordinate $\alpha\in\{u,v\}$ to source depth, where
$J^\alpha_{qk}\in\R$ has units of pixels per millimetre. The footprint radius
$r^\alpha_{qk}\in\R_+$, measured in pixels, is
\begin{equation}
 r^\alpha_{qk}
 =b+|J^\alpha_{qk}|\sigma_q+b_n(1-|n_q^z|).
\label{eq:dss-footprint}
\end{equation}

Here, $b,b_n\in\R_+$ are pixel-valued radii controlling the base footprint and
orientation-aware expansion, respectively. The scalar
$n_q^z\in[-1,1]$ is the optical-axis component of the orientation proxy. For
target pixel $p\in\Omega_t\subset\mathbb Z^2$, the dimensionless scalar
$\rho_{qkp}\in\R_+$ denotes its elliptical distance from the projected
hypothesis, normalized by both radii. The corresponding dimensionless
contribution $w_{qk\rightarrow p}\in\R_+$ is
\begin{equation}
 w_{qk\rightarrow p}=c_q a_k\exp\!\left[
 -\frac{\rho_{qkp}^2}{2}
 -\frac{[z^t_{qk}-z_{\min}(p)]_+}{\tau_z}\right].
\label{eq:dss-weight}
\end{equation}

Source confidence and quadrature mass scale this contribution, while spatial
distance and soft nearest-surface visibility form the exponential penalties.
Here, $z_{\min}(p)\in\R_+$ is the detached nearest arrival, $[\,\cdot\,]_+$ is
the positive part, and $\tau_z\in\R_+$ is a depth scale controlling
depth-order attenuation. Nearest-depth
selection and integer assignment are discrete; all remaining operations are
differentiable.

% \begin{figure}[t]
% \centering
% \includegraphics[width=\columnwidth]{figures/distrinvs_dss_evidence/distrinvs_dss_evidence.pdf}
% \caption{DSS evidence and reliability-guided expert routing for a held-out iMED
% validation target frame with substantial source-view disocclusion. (a) The
% completed-geometry DSS warp $\I_t^w$; unsupported pixels are black. (b)
% Support $S$, normalized only for display; the cyan contour marks the validity
% boundary. (c) The render-validity mask $M^{\rm dss}$ obtained from raw support.
% (d) The effective synthesis gate $g$ from
% Eq.~\eqref{eq:effective-gate}. It is zero on reliable support, selecting the
% bounded transport expert, and assigns greater weight to the
% synthesis expert as reliability decreases. The gate uses the full reliability
% condition rather than the complement of $M^{\rm dss}$.}
% \label{fig:dss-evidence}
% \vspace{-0.4cm}
% \end{figure}


\textbf{Reliability aggregation (RA).}
Let $\Omega_p\subseteq\Omega_s\times\mathcal K$ collect the retained
source--hypothesis contributions to target pixel $p$. Their total weight
defines raw support $S_p$. At render-valid pixels,
DSS normalizes these weights over $\Omega_p$. Their weighted moments yield the
transported tensor, expected depth, and conditional depth variance. Coverage
is a bounded saturation of raw support.
DSS couples normalized weight entropy to depth dispersion. Let
$H_N(p)\in[0,1]$ denote the normalized Shannon entropy of the retained weights
at pixel $p$. Collision
ambiguity is
\begin{equation}
 \mathcal H_p
 =H_N(p)\left[1-\exp(-V_p/\tau_z^2)\right].
\label{eq:collision-entropy}
\end{equation}

This coupling discounts entropy caused only by benign subpixel overlap. We
derive target confidence $\kappa_p$ by penalizing coverage with depth variance
and collision ambiguity. A pixel is render-valid when its raw support reaches
the dimensionless threshold $S_{\min}\in\R_+$.
% , as shown in
% Fig.~\ref{fig:dss-evidence}(c). Invalid pixels retain
% support and coverage but zero the transported outputs and ambiguity maps. The
% continuous evidence therefore distinguishes weak support from conflicting depth
% order rather than reducing reliability to a binary mask.

\subsection{Reliability-Guided Expert Fusion}

\textbf{Physics-guided routing (PGR).}
Both experts receive a condition tensor
$\mathbf h\in\R^{6\times H\times W}$, formed by concatenating coverage,
confidence, normalized variance, collision entropy, normalized depth, and DSS
validity. A fixed monotonic combination of these cues defines a synthesis prior
$g^0\in[0,1]^{1\times H\times W}$: low geometric reliability increases its
value, and pixels outside DSS
support are assigned entirely to synthesis. A residual router makes only a
bounded correction to the detached prior logit,
\begin{equation}
 \widetilde g_p=\operatorname{sigmoid}\!\left[
 \ell_p^0+\Delta_\theta(\mathbf h)_p\right],
 \qquad \ell^0=\operatorname{sg}[\operatorname{logit}(g^0)].
\label{eq:router}
\end{equation}
Here, $\widetilde g\in[0,1]^{1\times H\times W}$ is the learned gate,
$\ell^0\in\R^{1\times H\times W}$ is the detached prior logit, and
$\operatorname{sg}$ stops gradients. The bounded residual correction
$\Delta_\theta(\mathbf h)\in\R^{1\times H\times W}$ is predicted by the
router. A binary reliable-support mask
$M^{\rm rel}\in\{0,1\}^{1\times H\times W}$ accepts DSS evidence only when
the render is valid, sufficiently supported, and depth-consistent. The
effective synthesis gate $g\in[0,1]^{1\times H\times W}$ is
\begin{equation}
 g=(1-M^{\rm rel})\odot\widetilde g.
 \label{eq:effective-gate}
\end{equation}
Routing therefore retains a geometric initialization, suppresses synthesis on
reliable DSS support, and adapts elsewhere.
Fig.~\ref{fig:dss-evidence}(d) visualizes the resulting gate. It is not the
complement of $M^{\rm dss}$ because the continuous evidence and learned
residual also affect its value.

\textbf{Complementary experts.}
The transport expert predicts a bounded RGB correction to the DSS warp and
retains the DSS expected depth. The synthesis expert is an
uncertainty-conditioned Fourier U-Net~\cite{ronneberger2015unet} that combines
local convolution with real-FFT processing. The expert predictions satisfy
$\I_t^{\rm trans},\I_t^{\rm syn}\in[0,1]^{3\times H\times W}$,
$\D_t^{\rm syn}\in\R_+^{1\times H\times W}$, and
$\U_t^{\rm syn}\in[0,1]^{1\times H\times W}$.
This local--global design follows Fourier inpainting networks
~\cite{suvorov2022lama,chu2023uffc}, but all weights are trained within
DistriNVS. Risk is trained as an error-scale proxy on unreliable
regions, not as a calibrated failure probability over the full image.

\textbf{Reliability-constrained output.}
The transport branch retains source evidence while allowing bounded appearance
correction; the synthesis branch supplies content where geometric evidence is
weak. Their final predictions are
\begin{align}
 \widehat\I_t&=(1-g)\odot\I_t^{\rm trans}+g\odot\I_t^{\rm syn},\\
 \widehat\D_t&=(1-g)\odot\bar D_t+g\odot\D_t^{\rm syn},
 \qquad \U_t=\U_t^{\rm syn}.
 \label{eq:expert-fusion}
\end{align}
The gate is broadcast across RGB channels. We refer to this output rule as
\emph{soft fusion (SF)}. On $M^{\rm rel}$, $g=0$ and the output follows the
bounded transport expert; elsewhere the two experts form a soft mixture.
For comparison, \emph{hard composition (HC)} overrides the SF result on
$M^{\rm rel}$ by copying the unmodified DSS RGB $\I_t^w$ and depth $\bar D_t$.
Outside $M^{\rm rel}$, it retains the SF result rather than selecting the
synthesis expert alone. HC is evaluated only in the ablation study.

\subsection{Learning Objectives}

All loss terms defined below are nonnegative scalars in $\R_+$.
The training target is accompanied by a tissue-valid mask that excludes
annotated instruments when available. A regional reconstruction objective
$\mathcal L_{\rm rec}$ combines robust RGB and gradient agreement across trusted,
unreliable, and transition regions. These regions are determined by the
reliability masks. A depth objective
$\mathcal L_{\rm depth}$ combines a Laplace likelihood for the predicted source
depth distribution with relative target-depth supervision when target depth is
available.

For \emph{RGB-D-guided reprojection cycle consistency (RCC)}, the predicted
RGB-D is splatted back through the inverse calibrated transform and the
backward-rendered RGB is compared with the valid source image. Risk modulates
the backward splat, but no depth-cycle loss is imposed. The routing loss
$\mathcal L_{\rm router}$ uses binary cross-entropy against a detached
reliability target derived from DSS support, variance, and collision ambiguity.

The \emph{uncertainty loss (UL)} trains risk as an error scale on the
unreliable tissue mask $M_u\in\{0,1\}^{1\times H\times W}$. For pixel
$p\in\Omega_t$, let $e_p\in\R_+$ be the channel-mean absolute RGB error and
$s_p\in\R_+$ the positive scale predicted from risk. We use
\begin{equation}
 \mathcal L_{\rm UL}
 =\left\langle e_p/s_p+\log s_p\right\rangle_{M_u}.
\label{eq:risk-loss}
\end{equation}
The complete objective groups these terms as
\begin{equation}
 \mathcal L=\mathcal L_{\rm rec}+\lambda_d\mathcal L_{\rm depth}
 +\lambda_c\mathcal L_{\rm RCC}+\lambda_r\mathcal L_{\rm router}
 +\lambda_u\mathcal L_{\rm UL},
\label{eq:total-loss}
\end{equation}
where $\lambda_d,\lambda_c,\lambda_r,\lambda_u\in\R_+$ are scalar weights that
balance geometry, cycle, routing, and uncertainty supervision, respectively.

At inference, the predictor consumes only source RGB-D, camera calibration,
and relative pose, and ends at the fused prediction in
Eq.~\eqref{eq:expert-fusion}. Target RGB-D and instrument masks are used only
for supervision or evaluation; the backward-cycle branch and deterministic
image repair are disabled. The method assumes known
zero-skew pinhole intrinsics, a rigid inter-endoscope transform, and an
axis-aligned approximation to projected uncertainty; these assumptions delimit
its present scope.

\section{Experiments}
\label{sec:experiments}

\subsection{Experimental Setup}
\label{sec:implementation}

DistriNVS training uses the frozen list of 13 iMED Task~2 scenes supplied in
the project configuration~\cite{imed2026}. Splits are made at scene level, not
at frame level. Online depth dropout creates both isolated and contiguous
missing regions.

EndoVis/SCARED Datasets~8 and 9 are reserved for final zero-shot evaluation;
neither their images nor metrics are used for model selection. Pair inclusion
is determined by a renderer-only pass using the \emph{raw} valid source depth.
We retain 275 pairs whose raw DSS support exceeds 70\% of the image. Learned
depth completion, target RGB, and target tool masks cannot influence this
decision. Every method receives the same scene--frame identifiers, without
method-side or evaluator-side frame rejection.

We compare against seven endoscopic Gaussian reconstruction systems and four
feed-forward completion systems: LDM, LaMa, LCM-LoRA, and Hunyuan-DiT
~\cite{rombach2022ldm,suvorov2022lama,luo2023lcmlora,li2024hunyuandit}. Every
prediction is evaluated at the same resolution on the same target frames. 
We use the step-5,000 fine-tuned Hunyuan-DiT with the fixed
100-step DDPM protocol.


We report PSNR~\cite{huynhthu2008psnr}, SSIM~\cite{wang2004ssim}, and LPIPS~\cite{zhang2018lpips}.
H-prefixed metrics are evaluated
on the canonical raw-DSS hole intersected with non-tool tissue support on iMED.
All EndoVis pixels are evaluated because target tool masks are unavailable.
Scores are recomputed with a common evaluator, and degenerate predictions are
retained. The tables report frame-wise means and sample standard deviations,
rounded to two decimal places. Latency is measured at batch size one on an
RTX~4090 after warm-up, excluding loading, decoding, evaluation, and disk I/O.
The Hunyuan-DiT input mask controls synthesis but not the common H region.

\noindent\textbf{Implementation.}
Images are processed at $640\times512$. DSS uses the three Gaussian-weighted
hypotheses described in Sec.~\ref{sec:method}, and the reliability encoder
transports 16 feature channels. The synthesis expert has four resolutions and
base width 32, and the complete network contains $5.79$\,M trainable
parameters. We train for 10,000 AdamW steps with learning rate
$2\times10^{-4}$, weight decay $10^{-4}$, and effective batch size 12. Training
uses BF16 mixed precision and seed 6666. Renderer, routing, mask, and loss
constants are fixed across corresponding configurations and recorded in the
released configuration; none is tuned on EndoVis.

\subsection{In-Domain Comparison}

Tabs.~\ref{tab:validation-comparison} and
\ref{tab:validation-comparison-hole} compare all methods on the same 1,392
iMED frames using shared tissue and hole masks.
The Gaussian systems use dataset depth and serve as reconstruction references,
whereas DistriNVS is evaluated by its direct SF output without post-processing.
Bold and underlining mark the best and second-best means. 
On the iMED test set, DistriNVS achieved the highest PSNR, SSIM, and
H-PSNR while remaining competitive in LPIPS. Fig.~\ref{fig:imed-qualitative}
presents the corresponding qualitative evidence. Its examples are the 100th
retained pairs from the lexicographically first and last validation scenes, a
rule fixed before prediction inspection.

\begin{table}[t]
\centering
\caption{Full-image comparison on the seven-scene iMED test set
($N=1{,}392$).}
\label{tab:validation-comparison}
\scriptsize
\setlength{\tabcolsep}{0pt}
\begin{tabular*}{\columnwidth}{@{}l@{\extracolsep{\fill}}ccc@{}}
\toprule
Method & PSNR $\uparrow$ & SSIM $\uparrow$ & LPIPS $\downarrow$ \\
\midrule
Deform3DGS~\cite{yang2024deform3dgs} & $15.83\pm1.87$ & $0.33\pm0.06$ & $0.48\pm0.10$ \\
Endo-4DGS~\cite{huang2024endo4dgs} & $16.78\pm1.65$ & $0.36\pm0.04$ & $0.46\pm0.08$ \\
Free-SurGS~\cite{guo2024freesurgs} & $17.09\pm2.14$ & $\underline{0.43\pm0.07}$ & $0.51\pm0.07$ \\
SurgicalGS~\cite{chen2025surgicalgs} & $17.79\pm1.95$ & $0.38\pm0.07$ & $0.46\pm0.07$ \\
PR-ENDO~\cite{kaleta2025prendo} & $15.79\pm1.27$ & $0.31\pm0.07$ & $0.55\pm0.07$ \\
EndoGS~\cite{zhu2024endogs} & $12.52\pm6.22$ & $0.38\pm0.06$ & $0.66\pm0.20$ \\
EndoGaussian~\cite{liu2025endogaussian} & $14.38\pm4.26$ & $0.37\pm0.24$ & $0.77\pm0.12$ \\
LDM~\cite{rombach2022ldm} & $17.18\pm2.83$ & $0.34\pm0.06$ & $0.44\pm0.08$ \\
LaMa~\cite{suvorov2022lama} & $19.48\pm1.44$ & $0.43\pm0.11$ & $\mathbf{0.35\pm0.06}$ \\
LCM-LoRA~\cite{luo2023lcmlora} & $19.47\pm1.73$ & $0.42\pm0.12$ & $0.39\pm0.04$ \\
Hunyuan-DiT~\cite{li2024hunyuandit} & $\underline{19.52\pm1.51}$ & $0.42\pm0.12$ & $\underline{0.38\pm0.04}$ \\
\midrule
DistriNVS (SF) & $\mathbf{22.95\pm2.33}$ & $\mathbf{0.59\pm0.11}$ & $0.40\pm0.05$ \\
\bottomrule
\end{tabular*}
\end{table}

\begin{table}[t]
\centering
\caption{Hole-region comparison on the seven-scene iMED test set
($N=1{,}392$).}
\label{tab:validation-comparison-hole}
\scriptsize
\setlength{\tabcolsep}{0pt}
\begin{tabular*}{\columnwidth}{@{}l@{\extracolsep{\fill}}ccc@{}}
\toprule
Method & H-PSNR $\uparrow$ & H-SSIM $\uparrow$ & H-LPIPS $\downarrow$ \\
\midrule
Deform3DGS~\cite{yang2024deform3dgs} & $14.01\pm1.81$ & $0.31\pm0.07$ & $0.63\pm0.06$ \\
Endo-4DGS~\cite{huang2024endo4dgs} & $15.79\pm1.85$ & $0.42\pm0.14$ & $0.60\pm0.06$ \\
Free-SurGS~\cite{guo2024freesurgs} & $17.86\pm2.88$ & $0.56\pm0.05$ & $0.56\pm0.06$ \\
SurgicalGS~\cite{chen2025surgicalgs} & $17.63\pm2.87$ & $0.47\pm0.16$ & $0.59\pm0.06$ \\
PR-ENDO~\cite{kaleta2025prendo} & $11.92\pm1.71$ & $0.10\pm0.06$ & $0.74\pm0.08$ \\
EndoGS~\cite{zhu2024endogs} & $12.52\pm6.68$ & $0.49\pm0.12$ & $0.74\pm0.17$ \\
EndoGaussian~\cite{liu2025endogaussian} & $15.86\pm3.60$ & $0.48\pm0.31$ & $0.69\pm0.09$ \\
LDM~\cite{rombach2022ldm} & $14.90\pm3.48$ & $0.37\pm0.11$ & $0.64\pm0.13$ \\
LaMa~\cite{suvorov2022lama} & $21.40\pm2.32$ & $0.63\pm0.05$ & $0.42\pm0.08$ \\
LCM-LoRA~\cite{luo2023lcmlora} & $\underline{22.53\pm2.23}$ & $\underline{0.70\pm0.06}$ & $\underline{0.42\pm0.04}$ \\
Hunyuan-DiT~\cite{li2024hunyuandit} & $20.87\pm2.07$ & $0.64\pm0.08$ & $0.49\pm0.03$ \\
\midrule
DistriNVS (SF) & $\mathbf{27.95\pm2.91}$ & $\mathbf{0.80\pm0.07}$ & $\mathbf{0.41\pm0.05}$ \\
\bottomrule
\end{tabular*}
\end{table}

Fig.~\ref{fig:hole-ratio} stratifies reconstruction quality by the fraction of
target tissue unsupported by the source view.

\begin{figure}[t]
\centering
\includegraphics[width=\columnwidth]{figures/hole_ratio/hole_ratio_source_unobserved.pdf}
\caption{Performance stratified by source-unobserved ratio on the iMED test
set. Curves and shaded regions show the frame-wise mean and sample standard
deviation, respectively, for (a) PSNR and (b) LPIPS.}
\label{fig:hole-ratio}
\end{figure}



\begin{figure*}[t]
\centering
\includegraphics[width=\textwidth]{figures/table_comparisons/table1_qualitative_comparison.pdf}
\caption{Qualitative in-domain comparison on iMED. (a) Gaussian reconstruction
from source RGB. (b) Image completion from the DSS warp. White labels report
full-image PSNR (dB); teal borders denote DistriNVS.}
\label{fig:imed-qualitative}
\vspace{-0.4cm}
\end{figure*}

\subsection{Zero-Shot Comparison}

We compare the five feed-forward methods on the EndoVis
dataset, without EndoVis adaptation, or frame rejection.
Bold and underlining in Tab.~\ref{tab:zeroshot-breakdown} mark the best and
second-best unrounded means.
On EndoVis, DistriNVS achieved the highest PSNR and SSIM, the lowest LPIPS, and the
second-highest H-PSNR. It also had the second-lowest latency, while
LaMa achieved the highest H-PSNR and the lowest latency. The reported
DistriNVS latency includes depth reliability estimation, DSS, routing, both
experts, and output composition.
The slower inference of DistriNVS is partly due to the computational cost of DSS, 
whose distributional projection and aggregation operations add 
overhead beyond the synthesis network itself.
Fig.~\ref{fig:endovis-qualitative} shows the qualitative results of the zero-shot comparison on the EndoVis datasets.

\begin{table}[t]
\centering
\caption{Zero-shot comparison on EndoVis Datasets~8 and 9 ($N=275$).}
\label{tab:zeroshot-breakdown}
\scriptsize
\setlength{\tabcolsep}{0pt}

\begin{tabular*}{\columnwidth}{@{}l@{\extracolsep{\fill}}ccccc@{}}
\toprule
Method
& PSNR $\uparrow$
& SSIM $\uparrow$
& LPIPS $\downarrow$
& H-PSNR $\uparrow$
& Time (s) $\downarrow$ \\
\midrule

LDM
& $17.44{\pm}1.86$
& $0.33{\pm}0.14$
& $0.81{\pm}0.17$
& $12.54{\pm}2.44$
& $2.57{\pm}0.15$ \\

LaMa
& $\underline{22.72{\pm}2.10}$
& $\underline{0.63{\pm}0.14}$
& $\underline{0.32{\pm}0.07}$
& $\mathbf{21.43{\pm}2.71}$
& $\mathbf{0.04{\pm}0.00}$ \\

LCM-LoRA
& $14.26{\pm}0.94$
& $0.22{\pm}0.04$
& $1.00{\pm}0.08$
& $14.23{\pm}2.61$
& $0.40{\pm}0.04$ \\

Hunyuan-DiT
& $15.79{\pm}2.43$
& $0.31{\pm}0.12$
& $0.85{\pm}0.15$
& $13.24{\pm}3.21$
& $13.04{\pm}0.07$ \\

\midrule
DistriNVS
& $\mathbf{23.47{\pm}2.01}$
& $\mathbf{0.66{\pm}0.12}$
& $\mathbf{0.23{\pm}0.04}$
& $\underline{21.11{\pm}2.93}$
& $\underline{0.08{\pm}0.01}$ \\

\bottomrule
\end{tabular*}

\end{table}



\begin{figure*}[t]
\centering
\includegraphics[width=\textwidth]{figures/table_comparisons/table2_zeroshot_qualitative_comparison.pdf}
\caption{Qualitative zero-shot comparison on EndoVis Datasets~8 and 9 without
target-domain adaptation. Completion methods use the same DSS warp and mask.
White labels report full-image PSNR (dB); teal borders denote DistriNVS.}
\label{fig:endovis-qualitative}
\vspace{-0.4cm}
\end{figure*}

\subsection{Ablation Study}

Tab.~\ref{tab:ablation} reports eight single-seed configurations on the iMED test set.
The primary model uses SF, whereas the six component
variants retain legacy HC and are therefore interpreted only relative to the
HC reference. Bold and underlining mark the best and second-best means.
SF achieved the strongest overall reconstruction in Tab.~\ref{tab:ablation}.
Within the HC family, removing DD causes the clearest joint degradation,
whereas DC, OF, PGR, RCC, and UL show smaller or region-dependent changes.
Because SF and HC are trained separately, the HC variants should not be read
as direct removals from the primary SF model.

\begin{table}[t]
\centering
\caption{Ablation study on the seven-scene iMED test set ($N=1{,}392$).}
\label{tab:ablation}
\scriptsize
\setlength{\tabcolsep}{0pt}

\begin{tabular*}{\columnwidth}{@{}l@{\extracolsep{\fill}}rrrr@{}}
\toprule
Variant
& PSNR $\uparrow$
& SSIM $\uparrow$
& H-PSNR $\uparrow$
& H-SSIM $\uparrow$ \\
\midrule

\multicolumn{5}{@{}l}{\textit{Primary output rule and legacy reference}} \\
DistriNVS (SF)
& $\mathbf{22.95\pm2.33}$
& $\mathbf{0.59\pm0.11}$
& $\mathbf{27.95\pm2.91}$
& $\mathbf{0.80\pm0.07}$ \\

DistriNVS (HC)
& $\underline{20.77\pm1.91}$
& $\underline{0.49\pm0.13}$
& $27.56\pm2.73$
& $0.78\pm0.07$ \\

\addlinespace
\multicolumn{5}{@{}l}{\textit{Component variants under legacy hard composition}} \\

w/o DC
& $20.68\pm1.89$
& $0.49\pm0.13$
& $\underline{27.76\pm2.48}$
& $\underline{0.79\pm0.06}$ \\

w/o DD
& $20.67\pm1.91$
& $0.48\pm0.13$
& $27.19\pm2.74$
& $0.77\pm0.07$ \\

w/o OF
& $20.65\pm1.92$
& $0.48\pm0.13$
& $27.73\pm2.71$
& $0.78\pm0.06$ \\

w/o PGR
& $20.75\pm1.89$
& $0.49\pm0.13$
& $27.41\pm2.72$
& $0.77\pm0.07$ \\

w/o RCC
& $20.72\pm1.91$
& $0.49\pm0.13$
& $27.41\pm2.73$
& $0.78\pm0.07$ \\

w/o UL
& $20.70\pm1.89$
& $0.48\pm0.13$
& $27.02\pm2.60$
& $0.77\pm0.07$ \\

\bottomrule
\end{tabular*}

\end{table}



\section{Discussion}
\label{sec:discussion}

Our ablation indicates that the output rule has the largest effect, while DD
provides the most consistent component-level benefit within the HC family.
We observe smaller or region-dependent effects from the remaining components.
This pattern is consistent with HC preserving residual transport error, whereas
SF permits bounded appearance correction. Because SF and HC differ during both
training and inference, their comparison does not isolate the output operation
alone. Our ablation also uses a single seed, and the component variants are
trained under HC; their interactions with the primary SF model therefore remain
untested. We further note that the fixed EndoVis benchmark is frame-weighted
with an uneven sequence distribution, although all methods use the same pairs
without evaluator-side repair.

We observe a remaining gap to LaMa on several specific metrics, which may stem
from differences in model capacity and training objectives. LaMa uses
approximately 51M parameters and combines global Fourier processing with
adversarial, feature-matching, and perceptual objectives that favor
appearance-oriented completion. By contrast, our model uses only 5.79M
parameters while allocating capacity to geometry transport, reliability
estimation, routing, RGB-D prediction, and uncertainty modeling. Endoscopic
tissue often exhibits smooth color variation and low-frequency texture, which
may favor LaMa's larger appearance-completion network. This difference may
explain its lower LPIPS on iMED and its small H-PSNR advantage on EndoVis.
Our results nevertheless show that DistriNVS achieves better zero-shot PSNR,
SSIM, and LPIPS on EndoVis while using about 1/9 of the parameters.
Future work will explore increasing the capacity of our synthesis expert and 
introducing perceptual or adversarial supervision
while retaining the geometry-derived
reliability constraints that distinguish transported evidence from synthesized
content.

\section{Conclusion}
\label{sec:conclusion}

This paper introduces a reliability-aware formulation of cross-endoscope
novel-view synthesis through distributional surface splatting and constrained
transport--synthesis fusion. Projection support, depth variance, and collision
ambiguity determine where synthesis is permitted, while joint RGB-D prediction
provides content where transport is unreliable. The iMED and zero-shot EndoVis
results support SF as the primary output formulation and identify DD as the
most consistent component within the HC ablation family. Multi-seed studies,
matched SF component variants, and broader external datasets are needed for
stronger causal and generalization claims.

\bibliographystyle{IEEEtran}
\bibliography{references}

\end{document}
